\documentclass[]{article}
\usepackage{graphicx}
\usepackage{hyperref}
%opening
\title{Production Process for HPC/HMI code}
\author{Nyle Systems Control Group}

\begin{document}

\maketitle

\begin{abstract}
	This procedure configures a new stock HMI with the current HPC interface software. The provided USB provisioning media contains the application, required dependencies, and automated installer. No internet connection, additional hardware, or manual file copying is required.
\end{abstract}

\section{Turn on the HMI}
Providing power to the HMI should start the HMI. When it is fully loaded, you will see a desktop. The HMI may take approximately 30 seconds to boot initially.

\section{Insert the HPC HMI provisional media}
Insert the HPC HMI provisional media into any of the available USB type A connectors.The USB stick will auto mount. Select open,  then double click on the "stock-hmi" application. 

The application will run, and the HMI may reboot in the process. When it is complete, you will see the HPC HMI Graphical user interface. 

Remove your HPC HMI provisional media.

\section*{Creating Replacement HMI Provisioning Media}

If your HPC HMI provisioning media is lost, damaged, or otherwise unusable, create a replacement USB drive as follows:

\begin{enumerate}
	\item On your PC, go to:\par
	[
	\url{https://github.com/NyleWaterHeatingSystems/stock_hmi}
	]
	
	\item Select \textbf{Download ZIP}, then extract the downloaded ZIP file.
	
	\item Open the extracted folder and navigate to the \texttt{BOOKWORM} directory.
	
	\item Insert a blank USB drive and copy the \emph{contents} of the \texttt{BOOKWORM} directory to the top level of the USB drive.
	
	\item Safely eject the USB drive.
	
\end{enumerate}

The USB drive is now ready to provision an HPC HMI.


\end{document}
